% Responsable: Danela (1), Tahiny (2), Anthony (2) · Ronda 4 (domingo 15, hasta 15.00)
% Cinco conclusiones de mínimo 80 palabras: (1) afirmación, (2) evidencia con cita, (3) implicación para el PFC o la profesión.

% Chavarria Tahiny
\textbf{Conclusi\'on (Tahiny Chavarr\'ia, 1).} Distinguir los requisitos de sistema de los de software evita asignar al c\'odigo responsabilidades que corresponden al conjunto de hardware, software y personas. Las normas 29148, 15288 y 12207 describen una derivaci\'on por niveles, desde el stakeholder hasta el componente de software, que exige trazabilidad entre ellos \cite{iso29148,iso15288,ieee12207}. En el SGR-IA de El Balc\'on, el objetivo RN-01 de reducir mermas se traza hasta el pron\'ostico RF-02 y hasta su meta de calidad RNF-05, y as\'i el equipo sabe por qu\'e existe cada requisito y qu\'e debe revisar si cambia. Para la profesi\'on, esta trazabilidad es lo que permite gestionar el cambio sin perder el v\'inculo con las necesidades del cliente.

\textbf{Conclusi\'on (Tahiny Chavarr\'ia, 2).} Un modelo de calidad convierte adjetivos vagos en requisitos que se pueden probar. ISO/IEC 25010:2023 organiza la calidad del producto en nueve caracter\'isticas \cite{iso25010}, y Glinz advierte que, sin esa precisi\'on, los requisitos ``no funcionales'' se interpretan de forma inconsistente \cite{glinz2007}. En El Balc\'on, ``que sea r\'apido'' se transform\'o en RNF-01, con m\'etrica, umbral y condici\'on de carga, y se justific\'o por qu\'e mantenibilidad y flexibilidad quedan fuera de esta versi\'on. Para la profesi\'on, esto significa que la calidad se negocia con el cliente antes de programar y se verifica con pruebas, no con opiniones.


\textbf{Conclusión (Anthony Naranjo, 1).}
La definición del límite y del contexto del sistema permite comprender que los requisitos no surgen únicamente de las funciones que se desea programar, sino también de las personas, procesos y condiciones externas que influyen sobre la solución. El framework de Pohl considera el contexto como parte fundamental de la Ingeniería de Requisitos y propone perspectivas para analizarlo sistemáticamente \cite{pohl2025}. En El Balcón Restaurant, este análisis permitió diferenciar las funciones internas del SGR-IA de actores como el propietario, los usuarios, los clientes y servicios externos. Para
el PFC, esta delimitación reduce ambigüedades de alcance y proporciona una base más sólida para especificar requisitos verificables y mantener la coherencia del proyecto.

\textbf{Conclusión (Anthony Naranjo, 2).}
La aplicación de un proceso estructurado de Ingeniería de Requisitos permite transformar las necesidades de un negocio real en artefactos que pueden ser revisados, negociados, validados y gestionados durante el desarrollo. Pohl
organiza este trabajo mediante las actividades centrales de elicitación, documentación y negociación, apoyadas transversalmente por la validación y la gestión \cite{pohl2025}. En el caso de El Balcón Restaurant, estas
actividades permiten relacionar los problemas actuales de pedidos e inventario con requisitos concretos y con las funciones de inteligencia artificial propuestas. Para la profesión, aplicar este proceso ayuda a reducir decisiones basadas en supuestos y mantiene el desarrollo vinculado con las necesidades reales del cliente.
